%! Boxplots and Comparing Distributions — Unit 1, Lesson 1.6 — AP Practice
% EFFL, AP Statistics variant: AP-format practice, assigned but NOT scored — the cover's score
% column reads NA for this component. The graded practice is the `homework` component that
% follows it at the back of the packet.
% The four MC items cover the four ways this topic is tested: 1 is the width-means-count
% distractor, 2 is a five-number-summary/IQR reading item, 3 is mean-vs-median position under
% right skew (EK UNC-1.M.2), and 4 is the item whose correct answer is that the display CANNOT
% answer the question asked. Free-response (d) spirals to Lessons 1.3 and 1.4 (which display
% would show a second peak).
% Page lockstep with ap_practice_key rests on \writelines{n} in the blank matching a terse
% \ansline{} in the key — keep key answers short enough not to wrap past the reserved lines.
% The next-lesson preview is NOT here: it closes the packet, in ../homework/.
\documentclass[10pt]{article}
\usepackage{apstats-article}
\usepackage{apstats-boxes}
% \bxp{y-center}{half-height}{min}{Q1}{median}{Q3}{max} — the display is PRE-DRAWN.
% Authored byte-identically in the blank and the key.
\newcommand{\bxp}[7]{%
  \draw[thick, navy] (#3,#1) -- (#4,#1);
  \draw[thick, navy] (#6,#1) -- (#7,#1);
  \draw[thick, navy] (#3,{#1-#2}) -- (#3,{#1+#2});
  \draw[thick, navy] (#7,{#1-#2}) -- (#7,{#1+#2});
  \draw[thick, navy, fill=skymid] (#4,{#1-#2}) rectangle (#6,{#1+#2});
  \draw[very thick, navy] (#5,{#1-#2}) -- (#5,{#1+#2});}

\begin{document}

\pageheader{Unit 1, Lesson 1.6}{AP Practice}
% NAMESTRIP: no \namedateperiod here — the name row lives on the cover only.

\vspace{0.15in}

% Authored BYTE-IDENTICALLY in ap_practice/ and ap_practice_key/ — this box is what keeps the
% blank and the key the same number of pages.
\boxguard[8]
\begin{remindbox}
\small
This page is \textbf{not required and not scored} --- your graded homework is the last section
of this packet. These are AP-format reps: work them for the practice, and bring what you get
stuck on to office hours or study hall.
\end{remindbox}

\vspace{0.12in}

\boxguard[26]
\begin{scenariobox}[Weekday and Weekend Tips]{navy}
\small
A restaurant recorded one server's total tips, in dollars, on every weekday shift and every
weekend shift over three months. The two sets are shown below on a shared axis, with the
five-number summaries beneath.

\vspace{6pt}
\begin{center}
\begin{tikzpicture}[scale=0.85]
  \foreach \x in {0,1.4,2.8,4.2,5.6,7.0,8.4} {
    \draw[linegray, very thin] (\x,-0.15) -- (\x,2.05);}
  \bxp{1.5}{0.32}{0.14}{1.05}{1.54}{2.10}{3.08}
  \bxp{0.5}{0.32}{1.40}{2.66}{3.36}{4.62}{6.30}
  \fill[navy] (8.26,0.5) circle (0.085);
  \node[font=\bfseries\scriptsize, anchor=east] at (-0.22,1.5) {Weekday};
  \node[font=\bfseries\scriptsize, anchor=east] at (-0.22,0.5) {Weekend};
  \draw[->, thick] (-0.20,-0.15) -- (9.0,-0.15);
  \foreach \x/\lab in {0/40, 1.4/60, 2.8/80, 4.2/100, 5.6/120, 7.0/140, 8.4/160} {
    \draw (\x,-0.23) -- (\x,-0.07); \node[below, font=\scriptsize] at (\x,-0.27) {\lab};}
  \node[font=\scriptsize] at (4.2,-0.82) {Total tips for the shift (dollars)};
\end{tikzpicture}
\end{center}
\vspace{4pt}

\renewcommand{\arraystretch}{1.2}
\begin{center}
\begin{tabularx}{0.88\linewidth}{@{} l Y Y Y Y Y @{}}
\rowcolor{sky}
\textbf{Shift} & \textbf{Min} & \textbf{$Q_1$} & \textbf{Median} & \textbf{$Q_3$} &
\textbf{Max} \\
\hline
\textbf{Weekday} & \$42 & \$55 & \$62 & \$70 & \$84 \\
\textbf{Weekend} & \$60 & \$78 & \$88 & \$106 & \$158 \\
\end{tabularx}
\end{center}
\end{scenariobox}

\vspace{0.14in}

\boxguard[16]
\begin{headlinebox}{goldbg}
\small
\textbf{Section I --- Multiple Choice.} Circle the single best answer. Every answer choice is
about the context, not just the arithmetic --- read them all before choosing.
\end{headlinebox}

\vspace{0.10in}

\begin{enumerate}[leftmargin=1.9em, itemsep=0.13in]

  \item Which statement about the \textbf{weekend} shifts is best supported by the display?
        \begin{enumerate}[label=\textbf{(\Alph*)}, leftmargin=2.2em, itemsep=2pt]
          \item More weekend shifts brought in between \$106 and \$130 than between \$78 and
                \$88, because that stretch of the display is much wider.
          \item Half of the weekend shifts brought in between \$78 and \$88.
          \item The same number of weekend shifts fall between \$88 and \$106 as fall between
                \$78 and \$88.
          \item Weekend tips are spread evenly across the whole range from \$60 to \$158.
          \item Most weekend shifts brought in more than \$106.
        \end{enumerate}

  \item For the \textbf{weekday} shifts, the interquartile range of the daily tip totals is
        \begin{enumerate}[label=\textbf{(\Alph*)}, leftmargin=2.2em, itemsep=2pt]
          \item \$7
          \item \$15
          \item \$28
          \item \$42
          \item \$62
        \end{enumerate}

  \item The weekend distribution is skewed to the right and has one unusually high value. Which
        statement about the \textbf{weekend} tip totals is most likely correct?
        \begin{enumerate}[label=\textbf{(\Alph*)}, leftmargin=2.2em, itemsep=2pt]
          \item The mean is less than the median, because the long tail pulls the whisker out.
          \item The mean is approximately equal to the median, because the median is resistant.
          \item The mean is greater than the median, because the long right tail and the high
                value pull the mean upward.
          \item The mean is exactly \$88, the value marked inside the box.
          \item Nothing whatever can be said about the mean from these data.
        \end{enumerate}

  \item The manager asks whether weekend tip totals cluster around \emph{two} separate amounts
        --- a modest group of nights and a high group. What is the best response?
        \begin{enumerate}[label=\textbf{(\Alph*)}, leftmargin=2.2em, itemsep=2pt]
          \item Yes; the wide right side of the weekend display shows the second cluster.
          \item No; the display shows a single peak near \$88.
          \item Yes; the point plotted on its own is the second cluster.
          \item This display cannot answer the question --- it reports only five summary
                numbers, so clusters, gaps, and a second peak are invisible in it.
          \item No; tip data of this kind is always skewed right with one peak.
        \end{enumerate}

\end{enumerate}

\vspace{0.14in}

\boxguard[16]
\begin{headlinebox}{goldbg}
\small
\textbf{Section II --- Free Response.} Show your reasoning. Every answer must be written
\emph{in the context of the problem} --- that is what earns credit on the AP exam.
\end{headlinebox}

\vspace{0.10in}

\begin{enumerate}[leftmargin=1.9em, itemsep=0.16in, resume]

  \item Use the tip records above.

        \vspace{6pt}
        \textbf{(a)} Report the five-number summary for the \textbf{weekday} shifts, and say what
        each of the five numbers marks on the weekday display.\\[4pt]
        \writelines{2}

        \vspace{6pt}
        \textbf{(b)} For the \textbf{weekend} shifts, $Q_1 = \$78$ and $Q_3 = \$106$. Compute the
        interquartile range and use the $1.5 \times \text{IQR}$ rule to decide whether the \$158
        shift is an outlier. Show the arithmetic.\\[4pt]
        \writelines{3}

        \vspace{6pt}
        \textbf{(c)} Compare the weekday and weekend tip totals. Address center, variability, and
        unusual features, in context.\\[4pt]
        \writelines{4}

        \vspace{6pt}
        \textbf{(d)} The manager wants to know whether there are two distinct kinds of weekend
        night. Name a display you would ask for instead, and say what it could show that this one
        cannot.\\[4pt]
        \writelines{3}

\end{enumerate}

\vspace{0.12in}

\boxguard[12]
\begin{extensionbox}
\small
Find a pair of side-by-side boxplots published in a news article, a sports page, or a report.
Write one sentence comparing the two groups that uses a comparative word and names the variable
and its units --- then write the sentence the article actually used, and say which is stronger.\\[4pt]
\writelines{2}
\end{extensionbox}

\vspace{0.12in}


\end{document}
